%!TEX root = ../template.tex
%%%%%%%%%%%%%%%%%%%%%%%%%%%%%%%%%%%%%%%%%%%%%%%%%%%%%%%%%%%%%%%%%%%%
%% abstract-pt.tex
%% NOVA thesis document file
%%
%% Abstract in Portuguese
%%%%%%%%%%%%%%%%%%%%%%%%%%%%%%%%%%%%%%%%%%%%%%%%%%%%%%%%%%%%%%%%%%%%

\typeout{NT FILE abstract-pt.tex}%

\par O movimento \textit{Web3} tem promovido uma arquitetura mais descentralizada (e democrática) para a internet. Deste modo, sistemas descentralizados têm vindo a ganhar popularidade de novo. Um caso particular de sistemas \textit{peer-to-peer} aparece em \textit{Swarm Computing}, um paradigma em que dispositivos altamente heterogéneos cooperam diretamente de forma maioritariamente autónoma. Este paradigma exige soluções descentralizadas flexíveis e robustas, nomeadamente protocolos distribuidos e descentralizados de gestão de membros e comunicação.
\par A literatura \textit{peer-to-peer} propõe diversas abordagens para tais protocolos. No entanto, a grande maioria assume uma configuração estática, o que significa que os parâmetros que gerem a operação do sistema são definidos no momento do seu arranque e baseiam-se em expectativas sobre as condições do sistema ao longo do seu tempo de vida. Isto pode resultar num funcionamento ineficiente ou até na indisponibilidade do sistema. O paradigma da computação autónoma foi proposto há bastante tempo como um padrão de arquitetura para a construção de sistemas que se monitorizam a si próprios e ao seu ambiente para tomarem decisões sobre como se reconfigurarem, aplicando-as autonomamente. Soluções baseadas em computação autónoma tendem a ser altamente descentralizadas ou altamente centralizadas, o que abre espaço para o estudo de uma solução híbrida.
\par Neste trabalho desenhamos um componente centralizado (ou parcialmente centralizado) capaz de tomar decisões num sistema descentralizado em larga escala, o qual nomeámos \textit{Overlord}. Exploramos a ideia de computação autónoma para gerir parâmetros de protocolos \textit{peer-to-peer}. Para este fim, criamos o desenho da arquitetura, a sua implementação através da \textit{framework} Babel-Swarm, e a avaliação através de um caso de estudo de partilha de mensagens.

\keywords{
  Sistemas Descentralizados \and
  Computação Autónoma \and
  Peer-To-Peer \and
  Sistemas Auto-Adaptativos
}
% to add an extra black line
